% Resumo estruturado longo (Background/Methods/Results/Conclusions), com os trechos da
% revisao da Isa. NAO entra no manuscrito do IJF, que exige 100-150 palavras sem notacao.
% Guardado para revista que aceite resumo estruturado.

\section*{Abstract}

\noindent\textbf{Background.} Cardiovascular diseases are the leading cause of death in
Brazil. Foundation models pre-trained on large time series corpora are increasingly
reported as outperforming classical alternatives in health forecasting. Such rankings are
usually reported as point estimates, without any measure of how much of the observed gap
is sampling noise. We asked whether such a ranking survives a measure of its own
uncertainty, and where forecasting accuracy on this series actually comes from.

\noindent\textbf{Methods.} We benchmarked five forecasting approaches, SARIMA, Prophet,
TimesFM, XGBoost and CatBoost, on monthly cardiovascular mortality in the state of
S\~ao Paulo, using 168 months of records from the Brazilian Mortality Information System
(SIM/DataSUS, January 2010 to December 2023; 1{,}217{,}427 cardiovascular deaths).
Rolling origin cross-validation with a 6-month horizon and a 60-month minimum training
window produced 103 windows and 618 out-of-sample forecasts per model. Every comparison
carries uncertainty, from a block bootstrap resampling entire windows ($B=10{,}000$) and
from the Diebold-Mariano test with the Harvey correction at each horizon. Before running
the analysis we fixed the criterion for calling a difference established: the paired
interval must exclude zero and the Diebold-Mariano test must reach $p<0.05$ in at least
three of the six horizons. Three naive references were included as the benchmark any model
must clear. We further varied training history while holding test dates fixed, enriched
the feature set of the boosting models, added a tabular foundation model to the
comparison, and evaluated monthly minimum temperature as an exogenous covariate under a
leakage-free design.

\noindent\textbf{Results.} Prophet (sMAPE 4.70\%, 95\% CI 4.17 to 5.27), SARIMA (4.80\%,
4.29 to 5.33) and TimesFM (4.83\%, 4.32 to 5.39) were statistically indistinguishable:
the spread among them (0.13 percentage points, pp) was an order of magnitude smaller than
the mean interval width (1.07 pp), all three paired differences contained zero
($p=0.17$ to $0.77$), and none of the 18 Diebold-Mariano cells reached significance.
Boosting models were detectably worse, with all 36 comparisons against a leading model at
$p<0.05$, and neither cleared the seasonal naive benchmark (MASE 1.213 and 1.263 against
1.148); enriching their features and adding a tabular foundation model narrowed the gap
without closing it. Holding test dates fixed, nine additional years of training reduced sMAPE by
1.23 to 1.71 pp, roughly ten times any between-model difference. Truncating training
history to 60 months cost SARIMA 0.66 pp and Prophet 0.69 pp while leaving TimesFM
unaffected ($p=0.42$), and TimesFM then led. Minimum temperature improved three of the
four models that accept covariates, but only XGBoost met the pre-declared criterion (gain
0.368 pp, 95\% CI $+0.190$ to $+0.548$, 3 of 6 horizons), and a calendar term absorbed the
gain, indicating that what the covariate supplied was seasonality rather than weather.

\noindent\textbf{Conclusions.} Model selection among these three families is an
operational rather than a statistical decision. Extending the data delivered far more
than substituting the model. The defensible advantage of a zero-shot foundation model
here is robustness to short history, not superior accuracy.

